% TOP_PROBLEM: {"id":30007732,"problem_number":"LOCAL-30007732","title":"Persistence of early-childhood benefits","category_id":21,"category":{"id":21,"name":"economics","display_name":"Economics","description":"Open economic theory statements, published research agendas, and proposed scoped specifications; question provenance and historical priority are qualified per record.","slug":"economics","order_index":21,"created_at":"2026-10-08T00:00:00Z"},"status":"research_question","external_url":null,"legacy_ids":[],"published":false,"statement":"Estimate how later instructional improvements change preschool effects on skills through childhood and earnings through age twenty-five.","economics":{"original_id":221,"field":"Education, families and demography","kind":"empirical","formalization_basis":"adapted_research_question","source_status":"research_agenda","priority_status":"earliest_found","priority_note":"Earliest explicit statement located in this audit; no exhaustive priority search or first-ever claim. The mathematical specification below is adapted, not quoted from the source.","earliest_verified_reference":{"authors":["Greg J. Duncan","Ariel Kalil","Magne Mogstad","Mari Rege"],"title":"Investing in Early Childhood Development in Preschool and at Home","year":2022,"url":"https://bfi.uchicago.edu/wp-content/uploads/2022/05/BFI_WP_2022-58.pdf","doi":"10.3386/w29985","locator":"§7 Knowledge needs, printed pp. 91–92","evidence_summary":"The authors ask for long-term curriculum evidence and manipulation of follow-on school investments and combined preschool and parenting interventions."},"current_reference":{"authors":["Greg J. Duncan","Ariel Kalil","Magne Mogstad","Mari Rege"],"title":"Investing in Early Childhood Development in Preschool and at Home","year":2022,"url":"https://bfi.uchicago.edu/wp-content/uploads/2022/05/BFI_WP_2022-58.pdf","doi":"10.3386/w29985","locator":"§7 Knowledge needs, printed pp. 91–92","evidence_summary":"The authors ask for long-term curriculum evidence and manipulation of follow-on school investments and combined preschool and parenting interventions."},"formalization_note":"The source explicitly identifies the underlying gap or agenda. Population, horizon, interventions, estimands, and auxiliary assumptions are proposed operational choices, not a canonical source theorem. Partial later answers are compatible with this scoped question; the citation alone does not prove absence of every subsequent solution."},"statusQualification":"Published question or research agenda verified at the cited locator. The mathematical specification is adapted for this catalogue and includes additional assumptions; source publication does not establish first-ever priority or certify this exact model remains unsolved. The source explicitly identifies the underlying gap or agenda. Population, horizon, interventions, estimands, and auxiliary assumptions are proposed operational choices, not a canonical source theorem. Partial later answers are compatible with this scoped question; the citation alone does not prove absence of every subsequent solution.","statusReviewedAt":"2026-10-08"}
% FIRST_POSTED: 2026-10-08
% ENTRY_KIND: statement_only
% !TeX program = lualatex
\documentclass[11pt]{article}
\usepackage{fontspec}
\usepackage[a4paper,margin=25mm]{geometry}
\usepackage{amsmath,amssymb,mathtools}
\usepackage{xurl}
\usepackage[hidelinks]{hyperref}
\newcommand{\sref}[2]{\hyperlink{source:#1:#2}{[S#2]}}
\setlength{\emergencystretch}{3em}
\begin{document}
% problemId:30007732
\section*{Persistence of early-childhood benefits}\label{30007732}
\subsection{Definitions and mathematical statement}
Consider children eligible for a public preschool program in one country, followed from age four to age twenty-five. Let \(p\in\{0,1\}\) denote a randomized preschool offer and \(f\in\{0,1\}\) a separately randomized improvement in primary-school instruction beginning at age six. Let \(S_{ia}(p,f)\) be standardized skills at age \(a\), and \(Y_i(p,f)\) cumulative real earnings through age twenty-five. Define
\[\tau_a(f)=E[S_{ia}(1,f)-S_{ia}(0,f)],\qquad\tau_Y(f)=E[Y_i(1,f)-Y_i(0,f)].\]
The expectation covers baseline eligible children, including those who decline preschool, with outcome scales fixed independently of treatment. Identify the skill-response trajectory and earnings effects under factorial assignment, consistency, and a specified school-level spillover structure. Determine whether later instructional quality changes persistence of early gains, and whether convergence of measured test scores predicts convergence of later economic outcomes. Distinguish measurement artifacts, replacement of skills, and compensating parental investment where additional data permit. Report attrition and administrative-linkage uncertainty. The source explicitly requests long-term curriculum evidence and analysis of follow-on investments; this population, factorial variation, and age range make that agenda empirically reviewable rather than presume one universal fade-out mechanism.

\par\textbf{Formulation and scope.} The source explicitly identifies the underlying gap or agenda. Population, horizon, interventions, estimands, and auxiliary assumptions are proposed operational choices, not a canonical source theorem. Partial later answers are compatible with this scoped question; the citation alone does not prove absence of every subsequent solution.

\par\textbf{Open-question provenance.} An explicit question or research agenda was verified at the source locator. This does not by itself certify that every later version remains unresolved \sref{30007732}{1}.

\par\textbf{Historical priority.} Earliest explicit statement located in this audit; no exhaustive priority search or first-ever claim. The mathematical specification below is adapted, not quoted from the source.

\subsection{Short English statement}
Estimate how later instructional improvements change preschool effects on skills through childhood and earnings through age twenty-five.

\subsection{Sources}
\begin{itemize}\raggedright
\item[\textnormal{[S1]}]\hypertarget{source:30007732:1}{} Greg J. Duncan, Ariel Kalil, Magne Mogstad, Mari Rege. 2022. Investing in Early Childhood Development in Preschool and at Home. §7 Knowledge needs, printed pp. 91–92.
\url{https://bfi.uchicago.edu/wp-content/uploads/2022/05/BFI_WP_2022-58.pdf}.
DOI: 10.3386/w29985.
{\small\textit{Earliest verified question/agenda reference; historical priority qualified below. The authors ask for long-term curriculum evidence and manipulation of follow-on school investments and combined preschool and parenting interventions.}}
\end{itemize}
\end{document}
